
\section{Hướng dẫn sử dụng}

\subsection{Cách chạy chương trình}

\begin{enumerate}
    \item Mở MARS 4.5 simulator
    \item File → Open → Chọn file \texttt{battleship.asm}
    \item Run → Assemble
    \item Run → Go (hoặc nhấn F5)
\end{enumerate}

\subsection{Menu chính}

Khi chạy, chương trình hiển thị:

\begin{verbatim}
===== BATTLESHIP GAME =====
1. Play Game
2. Instructions
3. Quit
Choose: [nhập 1, 2, hoặc 3]
\end{verbatim}

\subsection{Chọn chế độ chơi}

Nếu chọn Play Game (1), chương trình yêu cầu:

\begin{verbatim}
Choose mode:
1. One Player (vs Computer)
2. Two Players
Choose: [nhập 1 hoặc 2]
\end{verbatim}

\subsection{Setup Phase}

Mỗi người chơi đặt 6 tàu. Nhập 4 số (hàng và cột của hai đầu tàu):

\begin{verbatim}
Player X: Place ship of length 2:
Enter: [nhập 4 số, ví dụ: 0 0 0 1]
\end{verbatim}

Nếu hợp lệ: thông báo "Ship placed". Nếu không: "Invalid ship. Try again."

\subsection{Battle Phase}

Mỗi lượt, nhập 2 số (hàng và cột):

\begin{verbatim}
Player X's turn:
[Hiển thị board]
Enter target: [nhập 2 số, ví dụ: 3 5]
\end{verbatim}

Kết quả: `HIT` hoặc `MISS`.

Hai Players mode: Sau mỗi lượt, chương trình tạm dừng để đổi lượt:

\begin{verbatim}
Press Enter to continue...
\end{verbatim}

\subsection{Ký hiệu trên board}

\begin{table}[H]
\centering
\begin{tabular}{|c|l|}
\hline
\textbf{Ký hiệu} & \textbf{Ý nghĩa} \\
\hline
`1` & Tàu còn sống (chỉ nhìn thấy trên board của mình) \\
\hline
`.` & Trống/chưa bắn \\
\hline
`O` & Trượt (MISS) \\
\hline
`X` & Trúng (HIT) \\
\hline
\end{tabular}
\caption{Ký hiệu trên board}
\end{table}

\subsection{One Player Mode (vs Computer)}

Computer sử dụng chiến lược đơn giản: quét từ index 0 đến 48, chọn ô đầu tiên chưa bắn. Không có AI phức tạp.


\begin{verbatim}
PROCEDURE game_loop()
  LOOP WHILE TRUE:
    
    +---> PLAYER 1 TURN
    |      |
    |      +---> CALL player_attack(Player1, Player2, memP1, boardP2)
    |             |
    |             +---> Display Player 1's board and memP1 (enemy map)
    |             +---> Read target row, column
    |             +---> Check if hit or miss
    |             +---> Update boardP2 and memP1
    |             +---> Return hit/miss result
    |
    |      +---> CALL board_has_ship(boardP2)
    |             |
    |             +---> IF no ship remains:
    |                    PRINT "Player 1 WINS!"
    |                    RETURN from game_loop
    |
    |      +---> IF Two Players mode:
    |             CALL pause_for_swap()
    |             (Display message and wait for confirmation)
    |
    +---> PLAYER 2 / COMPUTER TURN
    |      |
    |      +---> IF One Player mode:
    |             CALL computer_attack(Computer, Player1, memP2, boardP1)
    |      |     ELSE IF Two Players mode:
    |             CALL player_attack(Player2, Player1, memP2, boardP1)
    |             |
    |             +---> Display Player 2's board and memP2 (enemy map)
    |             +---> Read target row, column
    |             +---> Check if hit or miss
    |             +---> Update boardP1 and memP2
    |             +---> Return hit/miss result
    |
    |      +---> CALL board_has_ship(boardP1)
    |             |
    |             +---> IF no ship remains:
    |                    PRINT "Player 2 / Computer WINS!"
    |                    RETURN from game_loop
    |
    |      +---> IF Two Players mode:
    |             CALL pause_for_swap()

END LOOP
\end{verbatim}

\subsection{Chi tiết thủ tục player\_attack}

Thủ tục player\_attack xử lý một lượt tấn công của người chơi. Nó thực hiện các bước sau:

\subsection{Bước 1: Hiển thị bàn cờ}

In bàn cờ của người chơi hiện tại và bản đồ bắn đối phương (enemy map).

\begin{verbatim}
CALL print_dual_boards(own_board, enemy_memory_board)
\end{verbatim}

\subsection{Bước 2: Yêu cầu nhập tọa độ tấn công}

Yêu cầu người chơi nhập tọa độ hàng và cột mà người chơi muốn tấn công.

\begin{verbatim}
PRINT "Enter target row and column (e.g., '3 5'): "
READ input_line
CALL parse_two_ints(input_line, addr_row, addr_col)

IF parse failed:
  PRINT "Invalid input format. Please try again."
  RETURN
\end{verbatim}

\subsection{Bước 3: Kiểm tra phạm vi}

Kiểm tra tọa độ nhập vào có nằm trong [0, 6] hay không.

\begin{verbatim}
IF (row < 0 OR row > 6 OR col < 0 OR col > 6):
  PRINT "Coordinates out of range [0, 6]!"
  RETURN
\end{verbatim}

\subsection{Bước 4: Chuyển đổi tọa độ thành index}

Tính index từ tọa độ (row, column).

\begin{verbatim}
index = row * 7 + col
byte_offset = index * 4
\end{verbatim}

\subsection{Bước 5: Kiểm tra tấn công lặp lại}

Kiểm tra xem ô này đã được tấn công trước đó hay chưa bằng cách kiểm tra giá trị tương ứng trên bản đồ bắn (enemy\_memory\_board).

\begin{verbatim}
memory_value = enemy_memory_board[byte_offset]

IF memory_value != 0:  (memory_value là 1 hoặc 2)
  PRINT "You already attacked this cell! Try another."
  RETURN
\end{verbatim}

\subsection{Bước 6: Kiểm tra trúng hay hụt}

Kiểm tra giá trị tại ô đó trên bàn cờ của đối thủ.

\begin{verbatim}
opponent_value = opponent_board[byte_offset]

IF opponent_value == 1:
  // HIT
  PRINT "HIT!"
  
  // Update opponent's board
  opponent_board[byte_offset] = 0
  
  // Update attacker's memory
  enemy_memory_board[byte_offset] = 2
  
ELSE:  // opponent_value == 0
  // MISS
  PRINT "MISS!"
  
  // Update attacker's memory
  enemy_memory_board[byte_offset] = 1
\end{verbatim}

\subsection{Bước 7: Kiểm tra kết thúc trò chơi}

Sau mỗi lượt tấn công, kiểm tra xem bàn cờ của đối thủ còn tàu hay không.

\begin{verbatim}
CALL board_has_ship(opponent_board)

IF no ship remains:
  PRINT "You won! Game over."
  RETURN true (game over signal)
ELSE:
  RETURN false (game continues)
\end{verbatim}

\subsection{Thủ tục hỗ trợ: board\_has\_ship}

Thủ tục board\_has\_ship kiểm tra xem bàn cờ còn bất kỳ ô nào chứa tàu (giá trị = 1) hay không.

\begin{verbatim}
PROCEDURE board_has_ship(board)
  FOR i = 0 to 48:
    IF board[i] == 1:
      RETURN true  (still has ship)
  
  RETURN false  (no ship remains)

END PROCEDURE
\end{verbatim}

Nếu không còn ô nào có giá trị 1, thì tất cả tàu của người chơi đó đã bị phá hủy, và trò chơi kết thúc.

\subsection{Quá trình cập nhật bàn cờ}

Để hiểu rõ hơn cơ chế cập nhật dữ liệu, ta xem xét một ví dụ:

\subsection{Tình huống 1: Bắn trúng}

Giả sử Player 1 bắn vào (2, 3):
\begin{enumerate}
    \item index = 2 × 7 + 3 = 17
    \item byte\_offset = 17 × 4 = 68
    \item Kiểm tra memP1[17] (bản đồ bắn của Player 1):
    \begin{itemize}
        \item Nếu memP1[17] != 0: ô đã bắn rồi, từ chối.
        \item Nếu memP1[17] == 0: chưa bắn, tiếp tục.
    \end{itemize}
    \item Kiểm tra boardP2[17] (bàn cờ thực của Player 2):
    \begin{itemize}
        \item Nếu boardP2[17] == 1: trúng! Cập nhật boardP2[17] = 0, memP1[17] = 2.
        \item Nếu boardP2[17] == 0: trượt. Cập nhật memP1[17] = 1.
    \end{itemize}
\end{enumerate}

\subsection{Tình huống 2: Bắn hụt}

Nếu ô (2, 3) trên boardP2 là 0 (trống):
\begin{enumerate}
    \item Chỉ cập nhật memP1[17] = 1 (MISS).
    \item boardP2[17] vẫn là 0 (không thay đổi).
\end{enumerate}

\subsection{Tình huống 3: Bắn lặp lại}

Nếu memP1[17] đã là 1 hoặc 2 (đã bắn trước đó):
\begin{enumerate}
    \item Chương trình từ chối tấn công, yêu cầu nhập tọa độ mới.
    \item Không cập nhật bất kỳ dữ liệu nào.
\end{enumerate}

\subsection{Lưu ý quan trọng}

\begin{itemize}
    \item \textbf{Tách biệt boardP1/boardP2 và memP1/memP2}: Điều này rất quan trọng. Board chính chỉ được cập nhật khi đối thủ bắn trúng ô của mình. Memory board chỉ được cập nhật với kết quả tấn công của mình đối thủ.
    \item \textbf{Thứ tự kiểm tra}: Luôn kiểm tra memory board trước để đảm bảo không bắn lại cùng ô. Điều này giúp ngăn ngừa lỗi logic.
    \item \textbf{Hiển thị kết quả rõ ràng}: Sau mỗi lượt tấn công, in ra HIT! hoặc MISS! rõ ràng để người chơi biết kết quả.
\end{itemize}
